\chapter{Sensor Calibration}\label{app:a}

This appendix presents the raw calibration datasets used to derive the
correction equations for the voltage, current, and temperature sensing
channels. The data support the calibration procedures discussed in
\textit{Chapter~3} and the validation results presented in \textit{Chapter~4}.
The tables are retained in the appendix to show the measured sensor values
against the corresponding reference readings used during calibration.

\begin{table}[H]
  \centering
  \caption{Voltage Sensor Calibration Data}
  \label{tab:app-a-voltage-calibration}
  \footnotesize
  \setlength{\tabcolsep}{10pt}
  \renewcommand{\arraystretch}{0.95}
  \begin{tabular}{rrrr}
    \textbf{\shortstack{ZMPT101B\\(V)}} &
    \textbf{\shortstack{Fluke\\Multimeter (V)}} &
    \textbf{\shortstack{ZMPT101B\\(V)}} &
    \textbf{\shortstack{Fluke\\Multimeter (V)}} \\
    \hline
    34.1  & 35.2  & 179.1 & 180.5 \\
    43.9  & 44.79 & 189.0 & 190.0 \\
    58.6  & 60.1  & 199.2 & 200.4 \\
    73.0  & 74.5  & 204.7 & 205.3 \\
    83.3  & 85.0  & 210.7 & 210.0 \\
    99.3  & 100.2 & 214.2 & 214.8 \\
    113.6 & 115.1 & 220.0 & 220.2 \\
    129.0 & 130.0 & 225.3 & 225.1 \\
    139.6 & 140.3 & 228.4 & 230.8 \\
    149.0 & 150.0 & 235.3 & 235.1 \\
    159.1 & 160.0 & 238.8 & 240.7 \\
    169.2 & 169.9 & 242.9 & 244.6 \\
  \end{tabular}
\end{table}

The voltage calibration data show the uncorrected ZMPT101B readings and the
corresponding reference readings from the Fluke multimeter. These values were
used to generate the voltage correction equation applied in the firmware.

\begin{table}[H]
  \centering
  \caption{Current Sensor Calibration Data}
  \label{tab:app-a-current-calibration}
  \footnotesize
  \setlength{\tabcolsep}{9pt}
  \renewcommand{\arraystretch}{0.92}
  \begin{tabular}{rrrr}
    \textbf{\shortstack{ACS37035\\(A)}} &
    \textbf{\shortstack{Ingco Clamp\\Meter (A)}} &
    \textbf{\shortstack{ACS37035\\(A)}} &
    \textbf{\shortstack{Ingco Clamp\\Meter (A)}} \\
    \hline
    0.08 & 0.08 & 3.28 & 4.09 \\
    0.15 & 0.15 & 3.31 & 4.15 \\
    0.17 & 0.18 & 3.86 & 4.70 \\
    0.19 & 0.20 & 3.97 & 4.90 \\
    0.21 & 0.22 & 5.89 & 7.27 \\
    0.53 & 0.64 & 4.62 & 5.72 \\
    0.56 & 0.66 & 4.84 & 6.11 \\
    0.58 & 0.69 & 4.95 & 6.22 \\
    0.63 & 0.77 & 5.04 & 6.33 \\
    0.73 & 0.86 & 5.11 & 6.46 \\
    0.76 & 0.88 & 5.33 & 6.60 \\
    0.70 & 0.89 & 5.36 & 6.76 \\
    0.76 & 1.02 & 5.41 & 6.88 \\
    1.55 & 1.90 & 6.51 & 8.12 \\
    1.58 & 1.91 & 6.79 & 8.35 \\
    1.57 & 1.94 & 6.85 & 8.48 \\
    1.76 & 2.18 & 7.07 & 8.80 \\
    1.77 & 2.18 & 8.24 & 10.32 \\
    1.96 & 2.32 & 8.35 & 10.46 \\
    2.02 & 2.47 & 8.41 & 10.47 \\
    2.20 & 2.77 & 9.40 & 12.20 \\
    3.12 & 3.82 & 10.30 & 13.30 \\
    3.11 & 3.87 & 12.40 & 16.00 \\
  \end{tabular}
\end{table}

The current calibration data compare the ACS37035-derived current values with the
reference clamp-meter readings. These points were used to fit the current correction
equation used by the overload logic, socket-heating estimator, and waveform-feature
pipeline.

\begin{table}[H]
  \centering
  \caption{Reference Temperature and Thermistor Resistance Calibration Data}
  \label{tab:app-a-temperature-calibration}
  \footnotesize
  \setlength{\tabcolsep}{8pt}
  \renewcommand{\arraystretch}{0.95}
  \begin{tabular}{rrrr}
    \textbf{\shortstack{Reference\\Reading (\degreeC)}} &
    \textbf{\shortstack{Absolute\\Temperature (K)}} &
    \textbf{\shortstack{Measured Thermistor\\Resistance ($\Omega$)}} &
    \textbf{ln(R)} \\
    \hline
    38  & 311.15 K & 5.8k & 8.66072 \\
    50  & 323.15 K & 3.6k & 8.18869 \\
    100 & 373.15 K & 0.7k & 6.55108 \\
  \end{tabular}
\end{table}

The temperature calibration data were used to support the thermistor conversion
discussed in \textit{Chapter~3}. Since the thermistor measures the temperature at its
own placement point, these calibration values validate the sensor channel and do not
directly represent the hidden contact-spot temperature inside the socket interface.